\chapter*{Resumen}
\addcontentsline{toc}{chapter}{Resumen}
\markboth{Resumen}{Resumen}

\noindent Frente al software comercial de análisis por el Método de los Elementos Finitos (MEF), que opera como una caja negra, esta tesis desarrolla EduFEM, un software educativo libre y en español para la elasticidad lineal plana con cuadriláteros de cuatro y de nueve nodos (Q4 y Q9). La hipótesis sostiene que su desarrollo permitirá mejorar la forma en que el software expone el procedimiento de cálculo, elevando dos propiedades del análisis: que cada resultado pueda seguirse hasta sus datos (trazabilidad) y que cada magnitud pueda comprobarse contra un patrón independiente (verificabilidad). El método es el de la ciencia del diseño: se construye el software ---con un motor de cálculo, ocho módulos educativos y una memoria de cálculo automática cuyo desarrollo puede rehacerse a mano--- y se lo evalúa con criterios fijados de antemano. La trazabilidad se comprobó por cobertura: las nueve etapas del procedimiento quedan expuestas, y los dieciocho contenidos que un consenso de expertos exige dentro del alcance del trabajo tienen instrumento, según el cotejo del autor. La verificabilidad se comprobó con tres casos de estudio. Con soluciones manufacturadas ---campos exactos construidos para medir el error---, el error del desplazamiento se divide por cuatro con el Q4 y por ocho con el Q9 al reducir a la mitad el tamaño de los elementos, como predice la teoría. En la viga de Timoshenko, el error frente a la solución analítica fue del 0,04\,\% en la tensión normal $\sigma_x$ y del 0,26\,\% en la flecha, y la diferencia frente a SAP2000, del 0,21\,\% en $\sigma_x$. La membrana de Cook mostró el bloqueo por cortante del Q4 ---su rigidez excesiva bajo flexión--- frente a la convergencia del Q9. La hipótesis queda comprobada con esos criterios y dentro del alcance declarado: se evaluaron propiedades del software, no su uso por estudiantes, y los casos de contraste externo son de tensión plana.

\textit{Palabras clave:} método de elementos finitos, software educativo, trazabilidad, verificación y validación, memoria de cálculo
